% GENERATED FILE. Edit canonical lessons, then run npm run generate:books.
% canonical-source-hash: fnv1a64:e0d1b29ac3cbd4cf
% canonical-lessons: HI-C522-chach, HI-C522-mathna, HI-C522-jaman, HI-C522-khaulna, HI-C522-nitharna

\chapter{Milk and Curd}
\label{ch:hi-milk-and-curd}
\hlchaptermodality{522}

\begin{chapteropening}
\textbf{By the end of this chapter:} \emph{I can talk about churning, setting curd and draining whey.}

Complete the last lesson of chapter 522: I can talk about churning, setting curd and draining whey.
\end{chapteropening}

\section[\texorpdfstring{\hi{छाछ} (chāch)}{छाछ (chāch)}]{\hi{छाछ} (chāch) — buttermilk}

\label{lesson:HI-C522-chach}



\begin{quote}
Before the new one: say the Hindi for to temper (with hot spiced oil), then the Hindi for to roast.
\end{quote}

\subsection*{You'll want to know: \hi{छाछ}}
\textbf{\hi{छाछ}} — \emph{chāch} — \textquotedblleft{}buttermilk\textquotedblright{}.

\begin{quote}
\hi{गर्मी} \hi{में} \hi{हम} \hi{दोपहर} \hi{को} \hi{ठंडी} \hi{छाछ} \hi{पीते} \hi{हैं।}
\end{quote}

\begin{quote}
\emph{Garmī meṁ ham dopahar ko ṭhaṇḍī chāch pīte haiṁ.}
\end{quote}

\begin{quote}
In summer we drink cold buttermilk at midday.
\end{quote}

Feminine.

\begin{cousinweb}
\textbf{\hi{छाछ}} (\emph{chāch}) is a tadbhava, inherited from \emph{*chācchī}, \textquotedblleft{}buttermilk\textquotedblright{}, a form CDIAL reconstructs from its descendants (CDIAL 5012).
\end{cousinweb}

\begin{grammarlens}[title={What you've built}]
One more word for describing what you lived and felt.
\end{grammarlens}

\subsection*{Your turn}
\begin{itemize}
\raggedright
  \item \emph{Say it:} \emph{chāch}
  \item \emph{Say it:} \emph{chāch}, once more
  \item \emph{Say it:} \emph{Garmī meṁ ham dopahar ko ṭhaṇḍī chāch pīte haiṁ.} — In summer we drink cold buttermilk at midday.
  \item \emph{Recall:} say the Hindi for to temper (with hot spiced oil), then the Hindi for to roast, then say \emph{chāch} again
\end{itemize}

\begin{tcolorbox}[breakable,colback=teal!4,colframe=teal!35!black,title={Before you move on}]
What does \hi{छाछ} mean? (\textquotedblleft{}Buttermilk\textquotedblright{}.) Say it once more.
\end{tcolorbox}
\section[\texorpdfstring{\hi{मथना} (mathnā)}{मथना (mathnā)}]{\hi{मथना} (mathnā) — to churn}

\label{lesson:HI-C522-mathna}



\begin{quote}
Before the new one: say the Hindi for to roast, then the Hindi for buttermilk.
\end{quote}

\subsection*{You'll want to know: \hi{मथना}}
\textbf{\hi{मथना}} — \emph{mathnā} — \textquotedblleft{}to churn\textquotedblright{}.

\begin{quote}
\hi{दादी} \hi{दही} \hi{मथकर} \hi{मक्खन} \hi{निकालती} \hi{हैं।}
\end{quote}

\begin{quote}
\emph{Dādī dahī mathkar makkhan nikāltī haiṁ.}
\end{quote}

\begin{quote}
Grandmother churns yoghurt and takes out the butter.
\end{quote}

Past, with \hi{ने} (he or she …): \hi{उसने} \hi{दही} \hi{मथा}, \emph{usne dahī mathā}.

\begin{cousinweb}
\textbf{\hi{मथना}} (\emph{mathnā}) is a tadbhava, inherited from Sanskrit \emph{mathnāti}, \textquotedblleft{}stirs, churns\textquotedblright{} (CDIAL 9771).
\end{cousinweb}

\begin{grammarlens}[title={What you've built}]
One more word for describing what you lived and felt.
\end{grammarlens}

\subsection*{Your turn}
\begin{itemize}
\raggedright
  \item \emph{Say it:} \emph{mathnā}
  \item \emph{Say it:} \emph{mathnā}, once more
  \item \emph{Say it:} \emph{Dādī dahī mathkar makkhan nikāltī haiṁ.} — Grandmother churns yoghurt and takes out the butter.
  \item \emph{Recall:} say the Hindi for to roast, then the Hindi for buttermilk, then say \emph{mathnā} again
\end{itemize}

\begin{tcolorbox}[breakable,colback=teal!4,colframe=teal!35!black,title={Before you move on}]
What does \hi{मथना} mean? (\textquotedblleft{}To churn\textquotedblright{}.) Say it once more.
\end{tcolorbox}
\section[\texorpdfstring{\hi{जामन} (jāman)}{जामन (jāman)}]{\hi{जामन} (jāman) — a starter (for setting curd)}

\label{lesson:HI-C522-jaman}



\begin{quote}
Before the new one: say the Hindi for buttermilk, then the Hindi for to churn.
\end{quote}

\subsection*{You'll want to know: \hi{जामन}}
\textbf{\hi{जामन}} — \emph{jāman} — \textquotedblleft{}a starter (for setting curd)\textquotedblright{}.

\begin{quote}
\hi{रात} \hi{को} \hi{दूध} \hi{में} \hi{थोड़ा} \hi{जामन} \hi{मिलाइए}, \hi{सुबह} \hi{तक} \hi{दही} \hi{बन} \hi{जाएगा।}
\end{quote}

\begin{quote}
\emph{Rāt ko dūdh meṁ thoṛā jāman milāie, subah tak dahī ban jāegā.}
\end{quote}

\begin{quote}
Mix a little starter into the milk at night, and by morning it will be yoghurt.
\end{quote}

Masculine.

\begin{cousinweb}
\textbf{\hi{जामन}} (\emph{jāman}) is a tadbhava, inherited from Sanskrit \emph{yamyatē}, \textquotedblleft{}is held, is fixed\textquotedblright{} (CDIAL 10428).
\end{cousinweb}

\begin{grammarlens}[title={What you've built}]
One more word for describing what you lived and felt.
\end{grammarlens}

\subsection*{Your turn}
\begin{itemize}
\raggedright
  \item \emph{Say it:} \emph{jāman}
  \item \emph{Say it:} \emph{jāman}, once more
  \item \emph{Say it:} \emph{Rāt ko dūdh meṁ thoṛā jāman milāie, subah tak dahī ban jāegā.} — Mix a little starter into the milk at night, and by morning it will be yoghurt.
  \item \emph{Recall:} say the Hindi for buttermilk, then the Hindi for to churn, then say \emph{jāman} again
\end{itemize}

\begin{tcolorbox}[breakable,colback=teal!4,colframe=teal!35!black,title={Before you move on}]
What does \hi{जामन} mean? (\textquotedblleft{}A starter (for setting curd)\textquotedblright{}.) Say it once more.
\end{tcolorbox}
\section[\texorpdfstring{\hi{खौलना} (khaulnā)}{खौलना (khaulnā)}]{\hi{खौलना} (khaulnā) — to seethe; to boil hard}

\label{lesson:HI-C522-khaulna}



\begin{quote}
Before the new one: say the Hindi for to churn, then the Hindi for a starter (for setting curd).
\end{quote}

\subsection*{You'll want to know: \hi{खौलना}}
\textbf{\hi{खौलना}} — \emph{khaulnā} — \textquotedblleft{}to seethe; to boil hard\textquotedblright{}.

\begin{quote}
\hi{कड़ाही} \hi{में} \hi{तेल} \hi{खौल} \hi{रहा} \hi{है।}
\end{quote}

\begin{quote}
\emph{Kaṛāhī meṁ tel khaul rahā hai.}
\end{quote}

\begin{quote}
The oil in the wok is seething.
\end{quote}

Past (it …): \hi{वह} \hi{पानी} \hi{खौल} \hi{उठा}, \emph{vah pānī khaul uṭhā}.

\begin{cousinweb}
\textbf{\hi{खौलना}} (\emph{khaulnā}) is a tadbhava, inherited from \emph{*kha-ul}, \textquotedblleft{}be heated\textquotedblright{}, a form CDIAL reconstructs from its descendants (CDIAL 3760). CDIAL gives the Hindi verb as \textquotedblleft{}to be heated, boil\textquotedblright{}.
\end{cousinweb}

\begin{grammarlens}[title={What you've built}]
One more word for describing what you lived and felt.
\end{grammarlens}

\subsection*{Your turn}
\begin{itemize}
\raggedright
  \item \emph{Say it:} \emph{khaulnā}
  \item \emph{Say it:} \emph{khaulnā}, once more
  \item \emph{Say it:} \emph{Kaṛāhī meṁ tel khaul rahā hai.} — The oil in the wok is seething.
  \item \emph{Recall:} say the Hindi for to churn, then the Hindi for a starter (for setting curd), then say \emph{khaulnā} again
\end{itemize}

\begin{tcolorbox}[breakable,colback=teal!4,colframe=teal!35!black,title={Before you move on}]
What does \hi{खौलना} mean? (\textquotedblleft{}To seethe; to boil hard\textquotedblright{}.) Say it once more.
\end{tcolorbox}
\section[\texorpdfstring{\hi{निथारना} (nithārnā)}{निथारना (nithārnā)}]{\hi{निथारना} (nithārnā) — to drain off (liquid); to let settle and pour}

\label{lesson:HI-C522-nitharna}



\begin{quote}
Before the new one: say the Hindi for a starter (for setting curd), then the Hindi for to seethe.
\end{quote}

\subsection*{You'll want to know: \hi{निथारना}}
\textbf{\hi{निथारना}} — \emph{nithārnā} — \textquotedblleft{}to drain off (liquid); to let settle and pour\textquotedblright{}.

\begin{quote}
\hi{पनीर} \hi{बनाकर} \hi{माँ} \hi{ने} \hi{ऊपर} \hi{का} \hi{पानी} \hi{निथारा।}
\end{quote}

\begin{quote}
\emph{Panīr banākar māṁ ne ūpar kā pānī nithārā.}
\end{quote}

\begin{quote}
After making the paneer my mother drained off the water on top.
\end{quote}

Past, with \hi{ने} (he or she …): \hi{उसने} \hi{दाल} \hi{का} \hi{पानी} \hi{निथारा}, \emph{usne dāl kā pānī nithārā}.

\begin{cousinweb}
\textbf{\hi{निथारना}} (\emph{nithārnā}) is a tadbhava, inherited from Sanskrit \emph{nistārayati}, \textquotedblleft{}delivers from\textquotedblright{} (CDIAL 7532). CDIAL gives the Hindi verb as \textquotedblleft{}to draw off, purify\textquotedblright{}.
\end{cousinweb}

\begin{grammarlens}[title={What you've built}]
One more word for describing what you lived and felt.
\end{grammarlens}

\subsection*{Your turn}
\begin{itemize}
\raggedright
  \item \emph{Say it:} \emph{nithārnā}
  \item \emph{Say it:} \emph{nithārnā}, once more
  \item \emph{Say it:} \emph{Panīr banākar māṁ ne ūpar kā pānī nithārā.} — After making the paneer my mother drained off the water on top.
  \item \emph{Recall:} say the Hindi for a starter (for setting curd), then the Hindi for to seethe, then say \emph{nithārnā} again
\end{itemize}

\begin{tcolorbox}[breakable,colback=teal!4,colframe=teal!35!black,title={Before you move on}]
What does \hi{निथारना} mean? (\textquotedblleft{}To drain off (liquid); to let settle and pour\textquotedblright{}.) Say it once more.
\end{tcolorbox}
